\begin{frame}{diffusion이 여기에서 유리한 이유}
  \begin{itemize}
    \item 다중 에이전트 상호작용의 \textbf{결합 분포}를 배우므로,
      각 에이전트를 \textbf{독립 무작위}로 배치하는 것과 달리:
    \begin{itemize}
      \item ``\textbf{차선이 같은 차량들의 추종}''
      \item ``\textbf{보행자의 동시 이동}''
      같은 \textbf{상관 관계}를 잡아냄
    \end{itemize}
  \end{itemize}
  \vskip0.5em
  \begin{block}{조건화 $\Rightarrow$ ``조절 가능''}
    \textbf{자신의 상태와 지도}로 \textbf{조건화}하면,
    ``\textbf{그 교차로, 그 속도에서의 시나리오}''를 골라
    샘플링할 수 있음 --- ``\textit{조절 가능한 시뮬레이션
    초기화와 롤아웃}''이라는 논문 제목의 의미.
  \end{block}
  \vskip0.5em
  \begin{block}{역할}
    diffusion 모델이 \textbf{물리 시뮬레이터의 ``현실적인
    시작점''을 공급하는 생성적 사전분포(prior)}로 쓰이는 것 ---
    이것이 이 도구의 역할. 이후 도시 규모로 확장하는
    SceneDiffuser++(2025)도 ``\textbf{generative world
    model}''로 이어지고, 15.4 World Models와 \textbf{같은
    흐름}.
  \end{block}
\end{frame}
